% Part: sets-functions-relations
% Chapter: functions
% Section: inverses

\documentclass[../../../include/open-logic-section]{subfiles}

\begin{document}

\olfileid{sfr}{fun}{inv}
\olsection{فنکشناں دے اُلٹ}

\begin{explain}
اسی فنکشناں نوں نقشہ بندیاں سمجھدے آں۔ فیر فنکشناں بارے اک صاف
سوال ایہ اے کہ کی نقشہ بندی نوں ``اُلٹا'' چلایا جا سکدا اے۔ مثال
لئی، جانشین فنکشن $f(x) = x + 1$ نوں ایس معنی وچ اُلٹا چلایا جا
سکدا اے کہ فنکشن $g(y) = y - 1$، $f$ دے اثر نوں ``واپس موڑدا'' اے۔

پر سانوں احتیاط کرنا پئے گا۔ بھاویں~$g$ دی تعریف اک فنکشن
$\Int \to \Int$ تعریف کردی اے، پر ایہ اک \emph{فنکشن} $\Nat
\to \Nat$ تعریف نہیں کردی، کیوں جے $g(0) \notin \Nat$۔ سو سادیاں
صورتاں وچ وی ایہ بالکل صاف نہیں کہ فنکشن نوں اُلٹا چلایا جا سکدا اے
یا نہیں؛ ایہ دائرۂ تعریف تے ہدف سیٹ اُتے منحصر ہو سکدا اے۔

فنکشن دے اُلٹ دا خیال ایس گل نوں ہور ٹھیک ٹھیک بناؤندا اے۔
\end{explain}

\begin{defn}
اک فنکشن $g \colon B \to A$، فنکشن $f \colon A \to B$ دا
\emph{اُلٹ} اے جے $f(g(y)) = y$ تے $g(f(x)) = x$، ہر
$x \in A$ تے $y \in B$ لئی۔
\end{defn}

جے $f$ دا اُلٹ~$g$ ہووے، تاں اسی اکثر $f^{-1}$ لکھدے آں،~$g$ دی تھاں۔

\begin{explain}
ہُن اسی طے کراں گے کہ فنکشناں دے اُلٹ کدوں ہوندے نیں۔
$f\colon A \to B$ دے اُلٹ لئی اک چنگا امیدوار $g\colon B \to A$
اے جیہڑا ایویں ``تعریف کیتا'' جاوے:
\[
g(y) = \text{``اوہ'' $x$ جیہدے لئی $f(x) = y$۔}
\]
پر ``تعریف کیتا'' (تے ``اوہ'') دے گرد ڈراؤنے واوین دسدے نیں کہ ایہ
اجے تعریف نہیں۔ گھٹ توں گھٹ، ایہ پوری عمومیت نال ہمیشہ کم نہیں
کرے گی۔ ایس تعریف راہیں فنکشن ملن لئی اک، تے صرف اک،~$x$ ہونا
چاہیدا اے جیہدے لئی $f(x) = y$---یعنی~$g$ دا آؤٹ پٹ اکلوتے طور
تے طے ہونا چاہیدا اے۔ نالے ہر $y \in B$ لئی ایہ طے ہونا چاہیدا اے۔
جے $x_1$ تے $x_2 \in A$ ایہو جہے ہون کہ $x_1 \neq x_2$ پر
$f(x_1) = f(x_2)$، تاں $g(y)$، $y = f(x_1) = f(x_2)$ لئی اکلوتے
طور تے طے نہیں ہووے گا۔ تے جے کوئی~$x$ ای نہیں جیہدے لئی
$f(x) = y$، تاں $g(y)$ بالکل طے نہیں ہوندا۔ ہور لفظاں وچ،~$g$ دی
تعریف لئی~$f$ دا !!{injective} وی تے !!{surjective} وی ہونا ضروری اے۔

آؤ ہولی ہولی چلیے۔ اسی سوال نوں دو حصیاں وچ ونڈاں گے: دتا ہویا
فنکشن~$f\colon A \to B$ ہووے، تاں فنکشن $g\colon B \to A$ کدوں
ہوندا اے جیہدے لئی $g(f(x)) = x$؟ ایہو جہا~$g$، $f$ دے اثر نوں
``واپس موڑدا'' اے تے~$f$ دا \emph{کھبا اُلٹ} آکھلاندا اے۔ دوجا،
فنکشن $h\colon B \to A$ کدوں ہوندا اے جیہدے لئی $f(h(y)) = y$؟
ایہو جہا~$h$، $f$ دا \emph{سجا اُلٹ} آکھلاندا اے---$f$، $h$ دے
اثر نوں ``واپس موڑدا'' اے۔
\end{explain}

\begin{prop}
جے دائرۂ تعریف خالی نہیں تے $f\colon A \to B$، !!{injective} اے،
تاں اک \emph{کھبا اُلٹ}~$g\colon B \to A$،~$f$ دا اے جیہدے لئی
$g(f(x)) = x$، ہر $x \in A$ اُتے۔
% Source correction OLFUN-001: require nonempty A; exact condition is A nonempty or B empty.
\end{prop}

\begin{proof}
من لو کہ $f\colon A \to B$، !!{injective} اے۔ اک $y \in B$ اُتے
غور کرو۔ جے $y \in \ran{f}$، تاں کوئی $x \in A$ اے جیہدے لئی
$f(x) = y$۔ کیوں جے $f$، !!{injective} اے، ایس لئی ایہو جہا
اکو~$x \in A$ اے۔ فیر اسی تعریف کر سکدے آں: $g(y) = x$، یعنی
$g(y)$ ``اوہ'' $x \in A$ اے جیہدے لئی $f(x) = y$۔ جے
$y \notin \ran{f}$، تاں اسی ایہنوں کسے وی~$a \in A$ نال جوڑ سکدے
آں۔ دائرۂ تعریف خالی نہ ہون کرکے اسی اک $a \in A$ چُن سکدے آں تے
$g\colon B \to A$ نوں ایویں تعریف کر سکدے آں:
\[
g(y) = \begin{cases}
    x & \text{جے $f(x) = y$}\\
    a & \text{جے $y \notin \ran{f}$۔}
\end{cases}
\]
ایہ ہر $y \in B$ لئی تعریف شدہ اے، کیوں جے ایہو جہے ہر
$y \in \ran{f}$ لئی عین اک $x \in A$ اے جیہدے لئی $f(x) = y$۔
تعریف موجب، جے $y = f(x)$، تاں $g(y) = x$، یعنی
$g(f(x)) = x$۔
\end{proof}

\begin{prob}
وکھاؤ کہ جے $f\colon A \to B$ دا کھبا اُلٹ~$g$ اے، تاں~$f$،
!!{injective} اے۔
\end{prob}

\begin{prop}
    جے $f\colon A \to B$، !!{surjective} اے، تاں اک
    \emph{سجا اُلٹ}~$h\colon B \to A$،~$f$ دا اے جیہدے لئی
    $f(h(y)) = y$، ہر~$y \in B$ اُتے۔
\end{prop}

\begin{proof}
من لو کہ $f\colon A \to B$، !!{surjective} اے۔ اک $y \in B$ اُتے
غور کرو۔ کیوں جے~$f$، !!{surjective} اے، ایس لئی کوئی $x_y \in A$
اے جیہدے لئی $f(x_y) = y$۔ فیر اسی تعریف کر سکدے آں:
$h(y) = x_y$، یعنی ہر $y \in B$ لئی اسی کوئی $x \in A$ چُندے آں
جیہدے لئی $f(x) = y$؛ کیوں جے~$f$، !!{surjective} اے، ایس لئی
چُنن واسطے ہمیشہ گھٹ توں گھٹ اک ہوندا اے۔\footnote{کیوں جے $f$،
!!{surjective} اے، ہر~$y \in B$ لئی سیٹ $\Setabs{x}{f(x) = y}$
غیر خالی اے۔~$h$ دی ساڈی تعریف لئی سانوں ایہناں وچوں ہر سیٹ توں
اک $x$ چُننا پینا اے۔ ایہ گل کہ ہمیشہ ایویں چُنیا جا سکدا اے، اصل
وچ صاف نہیں---ایہ چوناں کر سکن دی صلاحیت بس اک مسلّم دے طور تے منی
گئی اے۔ ہور لفظاں وچ، ایہ قضیہ نام نہاد چُناؤ دے مسلّم نوں مندا اے؛
ایہ مسئلہ اسی اگے \oliflabeldef{sth:choice::chap}{
\olref[sth][choice][]{chap} وچ فیر چکاں گے}{نظر انداز کر دیواں گے}۔
پر کئی خاص صورتاں وچ، مثال لئی جدوں $A = \Nat$ ہووے یا محدود ہووے،
یا جدوں $f$، !!{bijective} ہووے، چُناؤ دے مسلّم دی لوڑ نہیں۔ (خاص
صورت وچ جدوں $f$، !!{bijective} ہووے، ہر $y \in B$ لئی سیٹ
$\Setabs{x}{f(x) = y}$ دا عین اک !!{element} ہوندا اے، ایس لئی
چُنن دی کوئی گل ای نہیں رہندی۔)}
تعریف موجب، جے $x = h(y)$، تاں $f(x) = y$، یعنی کسے وی
$y \in B$ لئی $f(h(y)) = y$۔
\end{proof}

\begin{prob}
وکھاؤ کہ جے $f\colon A \to B$ دا سجا اُلٹ~$h$ اے، تاں~$f$،
!!{surjective} اے۔
\end{prob}

\begin{explain}
  پچھلے ثبوت دے خیالاں نوں اکٹھا کرکے سانوں ہُن ایہ ملدا اے کہ ہر
  !!{bijection} دا اک اُلٹ ہوندا اے، یعنی اک اکو فنکشن جیہڑا~$f$
  دا کھبا اُلٹ وی اے تے سجا اُلٹ وی۔
\end{explain}

\begin{prop}\ollabel{prop:bijection-inverse}
جے $f\colon A \to B$، !!{bijective} اے، تاں اک
فنکشن~$f^{-1}\colon B \to A$ اے جیہدے لئی ہر $x \in A$ اُتے
$f^{-1}(f(x)) = x$ تے ہر $y \in B$ اُتے $f(f^{-1}(y)) = y$۔
\end{prop}

\begin{proof}
مشق۔
\end{proof}

\begin{prob}
\olref[sfr][fun][inv]{prop:bijection-inverse} ثابت کرو۔ تہانوں
~$f^{-1}$ تعریف کرنا اے، وکھاؤنا اے کہ ایہ فنکشن اے، تے وکھاؤنا
اے کہ ایہ~$f$ دا اُلٹ اے، یعنی $f^{-1}(f(x)) = x$ تے
$f(f^{-1}(y)) = y$، ہر $x \in A$ تے $y \in B$ لئی۔
\end{prob}

\begin{explain}
اُلٹ کڈھن دا اک ذرا ہور عام طریقہ وی اے۔ اسی \olref[kin]{sec} وچ
ویکھیا سی کہ ہر فنکشن $f$ اک !!a{surjection}
$f' \colon A \to \ran{f}$ پیدا کردا اے، $f'(x) = f(x)$ رکھ کے ہر
$x \in A$ لئی۔ صاف اے
کہ جے $f$، !!{injective} اے، تاں $f'$، !!{bijective} اے، سو
\olref{prop:bijection-inverse} موجب اوہدا اک اکلوتا اُلٹ اے۔ علامت
دی اک بہت چھوٹی جیہی ڈھیل نال اسی کدی کدی $f'$ دے اُلٹ نوں بس
``~$f$ دا اُلٹ'' آکھ دیندے آں۔
\end{explain}

\begin{prop}\ollabel{prop:left-right}%
  وکھاؤ کہ جے $f\colon A \to B$ دا کھبا اُلٹ~$g$ تے سجا
  اُلٹ~$h$ اے، تاں $h = g$۔
\end{prop}

\begin{proof}
  مشق۔
\end{proof}

\begin{prob}
  \olref[sfr][fun][inv]{prop:left-right} ثابت کرو۔
\end{prob}

\begin{prop}\ollabel{prop:inverse-unique}
ہر فنکشن~$f$ دا ودھ توں ودھ اک اُلٹ ہوندا اے۔
\end{prop}

\begin{proof}
  من لو کہ $g$ تے $h$ دونیں~$f$ دے اُلٹ نیں۔ فیر خاص طور تے
  $g$،~$f$ دا کھبا اُلٹ اے تے $h$ سجا اُلٹ۔
  \olref{prop:left-right} موجب، $g = h$۔
\end{proof}

\end{document}
